\item  \textbf{SAT Scores.}  A professor at a small univeristy in New York surveys her introductory statistics students at the beginning of each semester.  One of the questions she asks her students is what was the highest SAT score they received.  The professor has collected data over many semesters and has SAT scores for 353 students.  These data can be found on Blackboard (Assignments - Homework Assignments - Homework 2) in SAT.jmp/SAT.csv/SAT.sas/SAT.xlsx.  Use your favorite software to create a histogram and a boxplot of the SAT scores as well as to calculate summary statistics.  You do not need to give me the output, but you will need it to answer the following questions.
\begin{enumerate}
\item   Based on the output, describe the overall shape of the distribution of SAT scores for introductory statistics students that attended this small university in New York.
\item   Suppose you decide to use a Normal model to describe SAT scores for all introductory statistics students at this small university in New York.  What possible values of the parameters ($\mu$ and $\sigma$) could you use for this Normal model?
\item  The worst score a person can get on the SAT is a 400 and the best score a person can get is a 1600.  A Normal model allows for any value, including values less than 400 and greater than 1600.  Explain why it is still okay to use a Normal model to approximate the distribution of SAT scores.  
\end{enumerate}


